% Standalone insertion fragment. Requires amsmath, amssymb, and amsthm.
% All notation is defined below; no project-specific macros are used.
% This fragment does not assert a solution of the unrestricted coefficient problem.

\section{Variational selection at the topological defect scale}
\label{sec:o11-variational-selection-20260929}

We extract additional variational information from a sufficiently accurate
minimizing sequence. The argument uses finite metric variations directly;
it does not require a smooth minimizer, an Ekeland principle, or geometric
compactness. The test class is critical under four-dimensional scaling and
therefore admits variations supported on shrinking regions.

\paragraph{Setting and inherited estimates.}
Let
\[
 M_j=X_j\mathbin{\#}k_j\overline{\mathbb{CP}}^{\,2}
       \mathbin{\#}\ell_j(S^1\times S^3),
 \qquad A_j=c_1^2(X_j),\quad Q_j=k_j+4\ell_j>0,
 \quad h_j=Q_j/A_j\longrightarrow0,
\]
be the balanced complex-hyperbolic surgery family. For a smooth metric set
\[
 P(g)=\int_{M_j}|\operatorname{Ric}^{\circ}_g|^2\,dV_g,
 \qquad R(g)=\int_{M_j}|\operatorname{Ric}_g|^2\,dV_g,
 \qquad S(g)=\int_{M_j}s_g^2\,dV_g.
\]
Ricci norms are tensor norms. Curvature and Weyl norms below are the
Hilbert--Schmidt norms of their operators on two-forms. Define
\[
 p_j=\inf_{g\text{ on }M_j}P(g),\qquad
 \mathcal E(g)=4-
 \frac{\bigl(\int_{M_j}s_g\,dV_g\bigr)^2}
      {\operatorname{Vol}(g)R(g)},\qquad
 \alpha_j=\inf_{g\text{ on }M_j}\mathcal E(g).
\]
We use the following established properties of this family:
\begin{enumerate}
\item Every conformal class is negative Yamabe. Passing to its negative
Yamabe representative does not increase $P$, since
\[
 P(g)-P(\widehat g)=\frac{S(g)-S(\widehat g)}{12}\ge0.
\]
The representative can be scaled to $s_{\widehat g}=-4$ without changing
$P$.
\item The recovery construction gives $p_j\le8\pi^2Q_j$. For every metric,
\[
 8\pi^2A_j+P(g)\le R(g)
 \le8\pi^2A_jU_j+4P(g),\qquad U_j=1+O(h_j).
\]
\item The variational comparison gives
\[
 \alpha_j=\frac{4p_j}{8\pi^2A_j}+O(h_j^2).
\]
Moreover, $\chi(M_j)=A_j/3+O(Q_j)$ and
\[
 \int_{M_j}|\operatorname{Rm}_g|^2\,dV_g
       =8\pi^2\chi(M_j)+P(g).
\]
\end{enumerate}
All constants in these inherited estimates are uniform along the family.
The new argument below is conditional only on these displayed inputs.

\paragraph{Critical test budget.}
For a smooth symmetric covariant two-tensor $H$, put
\[
 \mathfrak D_g(H)=\int_M
      \bigl(|\nabla^2H|_g^2+|\nabla H|_g^4\bigr)\,dV_g.
\]
Under $(g,H)\mapsto(\lambda^2g,\lambda^2H)$, both
$\|H\|_{L^\infty(g)}$ and $\mathfrak D_g(H)$ remain unchanged.
In particular no pointwise bound on the derivatives of $H$ is imposed.

\paragraph{Lemma 1 (uniform finite-variation remainders).}
Suppose $g$ is a smooth metric on a closed four-manifold and
\[
 \int_M|\operatorname{Rm}_g|^2\,dV_g\le KA.
\]
If
\[
 \|H\|_{L^\infty(g)}\le1,\qquad \mathfrak D_g(H)\le A,
\]
then $g_t=g+tH$ is positive definite for $|t|\le t_0$, where $t_0>0$
is dimensional, and
\begin{equation}
 \left|\frac{d^m}{dt^m}P(g_t)\right|\le C_K A,
 \qquad m=2,3,\quad |t|\le t_0.
 \label{eq:o11-var-remainder-20260929}
\end{equation}
The constants do not require a diameter bound, injectivity-radius bound,
Sobolev inequality, or noncollapsing hypothesis.

\begin{proof}
The amplitude bound makes $g_t$ uniformly equivalent to $g$. In terms of
the background connection $\nabla=\nabla^g$, the difference tensor is
\[
 (\nabla^{g_t}-\nabla^g)^k{}_{ab}
 =\frac t2(g_t)^{kc}
   \bigl(\nabla_aH_{bc}+\nabla_bH_{ac}-\nabla_cH_{ab}\bigr).
\]
The curvature difference is the covariant derivative of this tensor plus
its quadratic product. Thus it consists of contractions of
\[
 t\,g_t^{-1}\nabla^2H
 \quad\text{and}\quad
 t^2g_t^{-1}g_t^{-1}\nabla H\nabla H.
\]
Changing contractions or lowering indices also involves the background
curvature and the uniformly equivalent metrics. Differentiating in $t$
through order three only inserts further uniformly bounded inverse
metrics and copies of $H$. Consequently the curvature, traceless Ricci
tensor, and their required $t$-derivatives are bounded pointwise by
\[
 C_m\bigl(|\operatorname{Rm}_g|+|\nabla^2H|+|\nabla H|^2\bigr).
\]
The volume density relative to $dV_g$, together with its first three
$t$-derivatives, is uniformly bounded. Differentiating
$|\operatorname{Ric}^{\circ}_{g_t}|_{g_t}^2dV_{g_t}$ and applying Young's
inequality to the cross terms bounds each integral in
\eqref{eq:o11-var-remainder-20260929} by
\[
 C\int_M\bigl(|\operatorname{Rm}_g|^2
                  +|\nabla^2H|^2+|\nabla H|^4\bigr)\,dV_g
 \le C_KA.
\]
No coordinate neighborhood of a uniform radius is used.
\end{proof}

\paragraph{Proposition 2 (weak stationarity and stability at scale $Q_j$).}
For every prescribed sequence $\varepsilon_j>0$ there are smooth metrics
$g_j$ on $M_j$ such that
\[
 s_{g_j}=-4,\qquad 0\le P(g_j)-p_j\le\varepsilon_j,
 \qquad P(g_j)-p_j\longrightarrow0,
 \qquad \mathcal E(g_j)-\alpha_j=o(h_j),
\]
and, uniformly over all smooth symmetric two-tensors satisfying
\[
 \|H\|_{L^\infty(g_j)}\le1,\qquad
 \mathfrak D_{g_j}(H)\le A_j,
\]
one has
\begin{equation}
 |DP_{g_j}[H]|=o(Q_j),\qquad
 D^2P_{g_j}[H,H]\ge-o(Q_j).
 \label{eq:o11-var-selection-20260929}
\end{equation}
Here $D^2P_{g_j}[H,H]$ means the second derivative along $g_j+tH$.
It is not being identified with an intrinsic Hessian at an assumed
critical point.

\begin{proof}
Choose any $\eta_j\downarrow0$ and set
\begin{equation}
 \delta_j=\min\left\{
       \varepsilon_j,j^{-1},\eta_j^3\frac{Q_j^3}{A_j^2}
                    \right\}.
 \label{eq:o11-var-delta-20260929}
\end{equation}
The definition of the infimum supplies a smooth metric with $P\le
p_j+\delta_j$. Replace it by the negative Yamabe representative and
scale to $s=-4$. The inherited conformal formula and scale invariance
preserve
\begin{equation}
             0\le P(g_j)-p_j\le\delta_j.
 \label{eq:o11-var-gap-20260929}
\end{equation}
Since $p_j\le8\pi^2Q_j$ and $\delta_j=o(Q_j)$, these metrics satisfy
$P(g_j)=O(Q_j)$. The two-sided Ricci estimate yields
\[
 R(g_j)=8\pi^2A_j+O(Q_j).
\]
At scalar curvature $-4$, $R=P+4\operatorname{Vol}$, and hence
\[
 \operatorname{Vol}(g_j)=2\pi^2A_j+O(Q_j).
\]
The curvature identity and Euler-characteristic estimate then give
\[
 \int_{M_j}|\operatorname{Rm}_{g_j}|^2\,dV_{g_j}=O(A_j).
\]
Thus Lemma 1 applies uniformly to the stated test class.

For such an $H$, put $F(t)=P(g_j+tH)$. Every $g_j+tH$ with
$|t|\le t_0$ is an admissible global competitor, so
\[
                  F(t)\ge p_j\ge F(0)-\delta_j.
\]
Taylor expansion for both signs of $t$, using $|F''|\le CA_j$, gives
\[
 |F'(0)|\le\frac{\delta_j}{t}+CA_jt,
 \qquad 0<t\le t_0.
\]
Choosing $t$ comparable to $\sqrt{\delta_j/A_j}$ gives
\begin{equation}
 |DP_{g_j}[H]|\le C\sqrt{A_j\delta_j}
       \le C\eta_j^{3/2}Q_j\sqrt{h_j}=o(Q_j).
 \label{eq:o11-var-first-20260929}
\end{equation}
To estimate the second variation, add the inequalities for $t$ and $-t$.
The bound $|F'''|\le CA_j$ gives
\[
 F''(0)\ge-\frac{2\delta_j}{t^2}-CA_jt.
\]
Taking $t$ comparable to $(\delta_j/A_j)^{1/3}$, which is at most a
constant times $\eta_jh_j$, yields
\begin{equation}
 D^2P_{g_j}[H,H]\ge-CA_j^{2/3}\delta_j^{1/3}
                        \ge-C\eta_jQ_j.
 \label{eq:o11-var-second-20260929}
\end{equation}
Because $\eta_j\to0$, this is the asserted lower bound $-o(Q_j)$.
Both estimates are uniform over the entire test class.

Finally, constant scalar curvature gives $\mathcal E=4P/R$. The Ricci
estimate implies
\[
 \mathcal E(g_j)=\frac{4P(g_j)}{8\pi^2A_j}+O(h_j^2).
\]
Comparing with the inherited formula for $\alpha_j$ and using
\eqref{eq:o11-var-gap-20260929} gives
\[
 0\le\mathcal E(g_j)-\alpha_j
       =O(h_j^2)+O(\delta_j/A_j)=o(h_j).
\]
The prescribed absolute accuracy and convergence to zero follow from
\eqref{eq:o11-var-delta-20260929}.
\end{proof}

The conclusion is a selection statement about sufficiently accurate
global near-minimizers. It does not say that every competitive sequence
with $P=O(Q)$ is approximately stationary or stable.

\paragraph{The resulting weak equation.}
Write $W_+(g)=\int|W_g^+|^2\,dV_g$ and $T_+=2\chi+3\tau$.
Chern--Weil gives
\[
                 P=4W_++S/12-8\pi^2T_+.
\]
At a closed constant-scalar-curvature metric with $s=-4$, the scalar
square variation is
\[
 DS_g[H]=8\int_M\langle\operatorname{Ric}^{\circ}_g,H\rangle_g\,dV_g.
\]
Indeed the general scalar variation, integrated against the constant
scalar curvature, has no surviving divergence terms; its gradient is
$-2s\operatorname{Ric}^{\circ}$.
Proposition 2 therefore yields
\begin{equation}
 4DW_{+,g_j}[H]
 +\frac23\int_{M_j}
       \langle\operatorname{Ric}^{\circ}_{g_j},H\rangle_{g_j}\,dV_{g_j}
       =o(Q_j)
 \label{eq:o11-var-weak-equation-20260929}
\end{equation}
uniformly in the critical test class. Keeping the Weyl first variation
in this form avoids importing a separate Bach-tensor sign or norm
convention.

\paragraph{Shrinking supports.}
For a test tensor of fixed amplitude on a region of radius $r$ with
controlled rescaled geometry, the estimates
\[
 |\nabla H|=O(r^{-1}),\qquad |\nabla^2H|=O(r^{-2}),\qquad
 \operatorname{Vol}(\operatorname{supp}H)=O(r^4)
\]
give $\mathfrak D_g(H)=O(1)$, even when $r\to0$.
This illustrates how a microscopic test can satisfy the budget. The
controlled rescaled geometry in this illustration is not a hypothesis
of Proposition 2. A fixed constant multiple of the displayed budget
can also be allowed by multiplying $H$ by a fixed sufficiently small
constant.

\paragraph{Corollary 3 (disjoint-region selection without smooth precompactness).}
Let $g_j$ be selected as in Proposition 2. For each $j$, let
$H_{ji}$ be finitely many smooth symmetric two-tensors with pairwise
disjoint supports, and suppose
\[
 \|H_{ji}\|_{L^\infty(g_j)}\le1,
 \qquad \sum_i\mathfrak D_{g_j}(H_{ji})\le A_j.
\]
Then
\begin{equation}
 \sum_i|DP_{g_j}[H_{ji}]|=o(Q_j),\qquad
 \sum_i\bigl(-D^2P_{g_j}[H_{ji},H_{ji}]\bigr)_+=o(Q_j).
 \label{eq:o11-var-multiregion-20260929}
\end{equation}

\begin{proof}
Choose signs $\epsilon_i\in\{-1,1\}$ to make
$DP[\epsilon_iH_{ji}]\ge0$, and put $H=\sum_i\epsilon_iH_{ji}$.
Disjointness makes the amplitude bound persist and the derivative
budgets add. The linearity of the first variation and Proposition 2 give
the first conclusion.

For the second conclusion, sum only the tensors whose second variations
are negative. The functional $P$ is the integral of a local expression
in the second jet of the metric. Mixed second variations therefore
vanish for disjoint supports. The second variation along the sum is
exactly the sum of the negative second variations. Proposition 2 again
gives the result.
\end{proof}

This corollary assumes neither a fixed smooth compact family of cell
metrics nor uniform collar geometry or cell scale. It does require the
actual test tensors to satisfy the total critical budget. In particular,
a family of defects admitting definite descent per unit defect cannot
have positive $Q_j$-density if its descent tests satisfy that budget.
The existence of such tests for all potentially inexpensive defects has
not been proved here.

\paragraph{An unresolved instability criterion for the sharp coefficient.}
The following is a sufficient criterion, not an established curvature
inequality. Suppose that for every fixed $\epsilon\in(0,1)$ there are
$c_\epsilon>0$ and $j_\epsilon$ such that every smooth metric on $M_j$,
$j\ge j_\epsilon$, satisfying
\[
 s_g=-4,\qquad P(g)\le8\pi^2(1-\epsilon)Q_j
\]
admits a smooth symmetric tensor with
\[
 \|H\|_{L^\infty(g)}\le1,\qquad \mathfrak D_g(H)\le A_j,
\]
for which at least one of
\begin{equation}
 |DP_g[H]|\ge c_\epsilon Q_j,
 \qquad
 D^2P_g[H,H]\le-c_\epsilon Q_j
 \label{eq:o11-var-open-instability-20260929}
\end{equation}
holds. Then
\[
 \frac{p_j}{8\pi^2Q_j}\longrightarrow1,
 \qquad \frac{\alpha_j}{h_j}\longrightarrow4.
\]

\begin{proof}[Proof of the conditional implication]
If the lower limit of $p_j/(8\pi^2Q_j)$ were less than one, choose a
subsequence with a fixed strict gap below one. The selected metrics of
Proposition 2 remain below a slightly weaker fixed threshold, because
$P(g_j)-p_j=o(Q_j)$. Both alternatives in
\eqref{eq:o11-var-open-instability-20260929} contradict the corresponding
uniform estimates in Proposition 2. Hence the lower limit is at least
one. The recovery upper bound gives the first limit, and the inherited
variational comparison gives the second.
\end{proof}

Proving \eqref{eq:o11-var-open-instability-20260929}, or constructing
genuine metrics with a smaller first-order cost, remains the substantive
unrestricted problem. In particular, one cannot insert
$H=\operatorname{Ric}^{\circ}$ in the preceding weak equation without
checking its amplitude and critical derivative budget. Nor do these
selection statements establish smooth convergence, a bubble-tree
decomposition, or the signed-curvature compensation inequality.
